\textbf{Topic 1: Data Representation\#2}\\
\subsection{Q\#1: Convert the following negative decimal integers into sign-and-magnitude and two’s complement representations (using 16 bits):}

\begin{tabularx}{\textwidth}{|c|X|X|}
\hline
\textbf{Decimal Number} & \textbf{Sign-and-Magnitude Representation} & \textbf{Two’s Complement Representation} \\
\hline -153 & 1000 0000 1001 1001 & 1111 1111 0110 0111 \\
\hline -421 & 1000 0001 1010 0101 & 1111 1110 0101 1011 \\
\hline -3276 & 1000 1100 1100 1100 & 1111 0011 0011 0100 \\
\hline -36317 & 
1000 1101 1101 1101 (Out of range, represents -3549) & 
0111 0010 0010 0011 (Out of range, represents 29219)\\
\hline
\end{tabularx}

\subsection{Q\#2: Convert the following binary numbers (using 16 bits) into decimal equivalent using sign-and-magnitude and two’s complement interpretations:}
\begin{tabularx}{\textwidth}{|c|c|c|}
\hline
\textbf{Binary Representation} & \textbf{Decimal using Sign-and-Magnitude} & \textbf{Decimal using Two’s Complement} \\
\hline
0000 0000 1011 1101 & 189 & 189 \\
\hline
1000 1100 1011 1001 & -3257 & -29511 \\
\hline
1000 1111 0000 1010 & -3850 & -28918 \\
\hline
1111 1111 1111 1111 & -32767 & -1 \\
\hline
\end{tabularx}
\subsection{Q\#3: Calculate the range of numbers that could be represented using different representations:}
\begin{table}[htbp]
\begin{tabularx}{\textwidth}{| X | X | X | X |}
\hline
\textbf{No of Bits} & \textbf{Unsigned Numbers(Min to Max)} & \textbf{Sign-and-Magnitude (Min to Max)} & \textbf{Two’s Complement (Min to Max)} \\
\hline
8 & 0 to 255 & -127 to +127 & -128 to +127 \\
\hline
16 & 0 to 65535 & -32767 to +32767 & -32768 to +32767 \\
\hline
32 & 0 to 4294967295 & -2147483647 to +2147483647 & -2147483648 to +2147483647 \\
\hline
64 & 0 to 1.844674407$\times$10\textsuperscript{19} & 
-9.223372037$\times$10\textsuperscript{18} to +9.223372037$\times$10\textsuperscript{18} & 
-9.223372037$\times$10\textsuperscript{18} to +9.223372037$\times$10\textsuperscript{18} \\
\hline
128 & 0 to 3.402823669$\times$10\textsuperscript{38} & 
-1.701411835$\times$10\textsuperscript{38} to +1.701411835$\times$10\textsuperscript{38} & 
-1.701411835$\times$10\textsuperscript{38} to +1.701411835$\times$10\textsuperscript{38} \\
\hline
\end{tabularx}
\end{table}

\subsection{Q\#4: What would be the result of adding a number to its one’s complement? What would be the result of adding an integer to its two’s complement?}
When you add a number to its one's complement, the result is always a bit pattern where all the bits are set to 1. The value of this number will depend on how we interpret it. For example, when using sign-and-magnitude representation (11111111\textsubscript{2}) will be equal to (-127)\textsubscript{10}, and in two’s complement representation it will be equal to (-1)\textsubscript{10}. In general, for an $n$-bit number, the result of adding a number to its one's complement will be $2^n - 1$.
When you add an integer to its two’s complement, the result is always zero. When you add the number $x$ and its two’s complement, the binary result overflows, leaving a sum of zero. This is because the two’s complement is effectively the negative representation of the original number in the binary system.
\subsection*{Q\#5: Normalize the following binary floating-point numbers. Explicitly show the value of the exponent after normalization.}
\textbf{a)} 1100.0111 \quad \textbf{b)} 10111.0011 × 2\textsuperscript{-5} \quad \textbf{c)} 0.0000010101 \quad \textbf{d)} 1101.001101 × 2\textsuperscript{7}

The normalized representations are given below:
\begin{itemize}
  \item a) 1.1000111 × 2\textsuperscript{3}
  \item b) 1.01110011 × 2\textsuperscript{1}
  \item c) 1.0101 × 2\textsuperscript{6}
  \item d) 1.101001101 × 2\textsuperscript{10}
\end{itemize}

\subsection{Q\#6: Convert the following decimal real numbers into fixed point representation, both in sign-and-magnitude and two’s complement representations (using 16 bits). For fixed point sign-and-magnitude, use 10 bits (including sign bit) for the integer part and 6 bits for the fractional part.}

\begin{tabular}{|c|c|c|}
\hline
\textbf{Decimal Number} & 
\begin{tabular}[c]{@{}c@{}}\textbf{Fixed Point Sign-and-Magnitude}\\ \textbf{Representation (10.6 format)}\end{tabular} & 
\begin{tabular}[c]{@{}c@{}}\textbf{Fixed Point Two’s Complement}\\ \textbf{Representation (10.6 format)}\end{tabular} \\
\hline
11.40625 & 0000001011.011010 & 0000001011.011010 \\
\hline
141.8125 & 0010001101.110100 & 0010001101.110100 \\
\hline
-412.21875 & 1110011100.001110 & 1001100011.110010 \\
\hline
-465.09375 & 1111010001.000110 & 1000101110.111010 \\
\hline
\end{tabular}

\textbf{Note:} For 2’s complement representation of a negative real number, we initially assume the number to be positive, convert to binary in the same way as fixed point conversion and apply two’s complement on the binary representation to get the final answer.

\subsection*{Q\#7:Convert the following decimal real numbers into fixed point representation (using 8 bits). Convert the fixed point back to decimal real numbers. Do you get the same value as the original?}
\begin{center}
\begin{tabularx}{\textwidth}{|X|X|X|X|}
\hline
\textbf{Decimal Number (Original)} & \textbf{Fixed Point Sign-and-Magnitude Representation (3.5 format)} & \textbf{Decimal Number (After Conversion)} & \textbf{Same as Original?} \\
\hline
3.635 & 011.10100 & 3.625 & No \\
\hline
-2.10 & 110.00011 & -2.09375 & No \\
\hline
-3.33 & 111.01010 & -3.3125 & No \\
\hline
3.14159 & 011.00100 & 3.125 & No \\
\hline
\end{tabularx}
\end{center}

\subsection{Q\#8: Convert the following numbers into IEEE 754 (32-bit) floating point format with 8-bit exponent:}
\begin{tabularx}{\textwidth}{|c|X|}
\hline
\textbf{Decimal Number} & \textbf{IEEE 754 Floating Point Representation (32-bit)} \\
\hline
\textbf{104.1875} &
Conversion to binary gives: 1101000.0011 \newline
Normalized representation: 1.1010000011 × 2\textsuperscript{6} \newline
Biased Exponent: 6 + 127 = 133 = (10000101)\textsubscript{2} \newline
Mantissa to be stored: (.1010000011 0000 0000 0000)\textsubscript{2} \newline
Final Representation: \newline
\texttt{0 10000101 10100000110000000000000} \\
\hline
\textbf{-12.640625} &
Conversion to binary gives: 1100.101001 \newline
Normalized representation: 1.100101001 × 2\textsuperscript{3} \newline
Biased Exponent: 3 + 127 = 130 = (10000010)\textsubscript{2} \newline
Mantissa to be stored: (.10010100100000000000000)\textsubscript{2} \newline
Final Representation: \newline
\texttt{1 10000010 10010100100000000000000} \\
\hline
\textbf{1125.10} &
Conversion to binary gives: 10001100101.000110011 \newline
Normalized representation: 1.00011001010000110011 × 2\textsuperscript{10} \newline
Biased Exponent: 10 + 127 = 137 = (10001001)\textsubscript{2} \newline
Mantissa to be stored: (.00011001010000110011)\textsubscript{2} \newline
Final Representation: \newline
\texttt{0 10001001 00011001010000110011} \\
\hline
\textbf{-0.0033} &
Conversion to binary gives: 0.00000000001110100000001000011\ldots \newline
Normalized representation: 1.10110000001000011\ldots × 2\textsuperscript{-9} \newline
Biased Exponent: -9 + 127 = 118 = (01110110)\textsubscript{2} \newline
Mantissa to be stored: (.10110000001000011010000)\textsubscript{2} \newline
Final Representation: \newline
\begin{tabular}{|c|c|c|}
\hline
1 & \texttt{0111 0110} & \texttt{1011 0000 0010 0001 1010 000} \\
\end{tabular}
\\
\hline
\end{tabularx}

\subsection{Q\#9: Given the following numbers in IEEE 754 (32-bit) floating point format with 8-bit exponent, calculate the decimal real numbers.}
\begin{longtable}{|c|p{10cm}|}
\hline
\textbf{IEEE 754 Floating Point Representation (32-bit)} & \textbf{Decimal Number} \\
\hline
0\ 1000\ 0111\ 101\ 1100\ 1100\ 0000\ 0000\ 0000 & 
Sign bit: 0 (+ve) \newline
Biased Exponent: (10000111)\textsubscript{2} = 135 \newline
Real Exponent: 135 – 127 = 8 \newline
Stored Mantissa: (.10111001100000000000000)\textsubscript{2} \newline
Actual Mantissa: (1.10111001100000000000000)\textsubscript{2} \newline
Normalized representation: 1.101110011 × 2\textsuperscript{8} \newline
Binary number: 11011100.11 \newline
Decimal value: \textbf{441.5} \\
\hline
0\ 0111\ 1000\ 011\ 0000\ 0000\ 0000\ 0000\ 0000 & 
Sign bit: 0 (+ve) \newline
Biased Exponent: (01111000)\textsubscript{2} = 120 \newline
Real Exponent: 120 – 127 = -7 \newline
Stored Mantissa: (.01100000000000000000000)\textsubscript{2} \newline
Actual Mantissa: (1.01100000000000000000000)\textsubscript{2} \newline
Normalized representation: 1.011 × 2\textsuperscript{-7} \newline
Binary number: 0.0000001011 \newline
Decimal value: \textbf{0.010742188} \\
\hline
1\ 1000\ 1010\ 101\ 0101\ 0000\ 0100\ 0000\ 0000 & 
Sign bit: 1 (-ve) \newline
Biased Exponent: (10001010)\textsubscript{2} = 138 \newline
Real Exponent: 138 – 127 = 11 \newline
Stored Mantissa: (.10101010000010000000000)\textsubscript{2} \newline
Actual Mantissa: (1.10101010000010000000000)\textsubscript{2} \newline
Normalized representation: 1.1010101000001 × 2\textsuperscript{11} \newline
Binary number: 11010101000001.00 \newline
Decimal value: \textbf{-3408.25} \\
\hline
1\ 0111\ 1010\ 100\ 0000\ 0000\ 0000\ 0000\ 0000 & 
Sign bit: 1 (-ve) \newline
Biased Exponent: (01111010)\textsubscript{2} = 122 \newline
Real Exponent: 122 – 127 = -5 \newline
Stored Mantissa: (.10000000000000000000000)\textsubscript{2} \newline
Actual Mantissa: (1.10000000000000000000000)\textsubscript{2} \newline
Normalized representation: 1.1 × 2\textsuperscript{-5} \newline
Binary number: 0.000011 \newline
Decimal value: \textbf{-0.046875} \\
\hline
\end{longtable}

\subsection{Q\#10:}
The data-type \textbf{Float} (32-bit total = 23-bit mantissa, 8-bit exponent) gives us about 7–8 significant decimal digits. \\
The data-type \textbf{Double} (64-bit total = 52-bit mantissa, 11-bit exponent) gives us about 15–16 significant decimal digits.

\vspace{0.5em}
Let's create a new data-type called \textbf{Triple}, which has 96-bits total; a 78-bit mantissa and a 17-bit exponent. Roughly how many significant decimal digits will this type give us?

\textcolor{blue}{
The mantissa has 78-bits, so (when we include the hidden bit) we get 79 significant bits; or $2^{79}$ values.
$2^{79} = 6.044629098 \times 10^{23}$ 
(we can work this out through various methods, e.g. using logarithms of base 10). Therefore we have
roughly 23–24 significant decimal digits.
}